% p2-04-e8-toda

\section{P2 §4. The E8/Toda Anchor}

4. The E8/Toda Anchor: Exact $\varphi$ Occurrences in Field
  Theory

     Figure (fig-p2-ch4-e8-toda). E8 / Toda triptych --- E8 mass spectrum / x squared minus x minus
                                         one root / golden gears.

  4.1 Zamolodchikov's E8 Integrable Field Theory

  In 1989, Zamolodchikov considered the critical Ising model --- a 2D conformal field
  theory with central charge c = 1/2 --- perturbed by a magnetic field operator sigma
  (Adv. Stud. Pure Math. 19 (1989) 641):

  A = ACFT + h $\in$t d2x sigma(x). 9

  The operator sigma has conformal dimension Deltasigma = 1/16 in each chiral sector
  (i.e., holomorphic dimension hsigma = hsigma = 1/16, so scaling dimension
  2Deltasigma = 1/8). Despite breaking conformal symmetry, the perturbed theory
  retains an infinite tower of conserved charges at spins

  s $\in$ {1, 7, 11, 13, 17, 19, 23, 29, …},

  which are exactly the exponents of the E8 Lie algebra modulo the Coxeter number h =
  30. The theory is integrable: scattering is purely elastic, factorized, and described by an
  exact S-matrix (Int. J. Mod. Phys. A 4 (1989) 4235). The particle spectrum contains
  exactly 8 species, with masses proportional to the components of the Perron-Frobenius
  eigenvector of the E8 Cartan matrix.

  [L1 --- Theorem] The complete mass spectrum is:

  The fundamental $\varphi$ relation is

  (m2)/(m1) = 2cos((pi)/(5)) = $\varphi$ = (1+5)/(2). 10

  This is an exact result. All eight masses lie in the splitting field of the cyclotomic
  polynomial Phi30(x) over Q, which contains Q(5) as a subfield.

  4.2 Experimental Confirmation in CoNb2O6
  [L2 --- Experimental] Coldea et al. (2010) measured the spin excitation spectrum of the
  quasi-one-dimensional Ising ferromagnet cobalt niobate (CoNb2O6) in a transverse
  magnetic field tuned to the quantum critical point (Science 327 (2010) 177, DOI:
  10.1126/science.1180085). Neutron scattering revealed two sharp modes at low
  energies, with the ratio of their frequencies approaching the golden mean to within
  experimental precision --- the specific prediction for the first two E8 meson particles.
  This is the first and most direct experimental realization of an E8-structured mass
  spectrum in a condensed matter system.

  Critical caveat [L6 --- Excluded inference]. The E8 symmetry here is an emergent
  effective symmetry of a near-critical 1D quantum magnet. It is not a fundamental gauge
  symmetry of high-energy physics. The appearance of $\varphi$ is a consequence of the specific
  algebraic structure of the E8 Cartan matrix within that effective theory. No inference
  about quark masses, gauge coupling constants, or any Standard Model parameter is
  supported by this result. The universality classes of 2D integrable condensed matter
  and 4D gauge theory are disjoint in current theory.

  4.3 Toda Field Theories and Non-Crystallographic Coxeter Groups
  Affine Toda field theories are integrable 2D quantum field theories defined by a
  Lagrangian of the form

  L = (1)/(2)$\partial$mu phia $\partial$mu phia - (m2)/(beta2) sumi=0r ni exp(beta alphai $\cdot$ phi), 11

  where alphai are the simple roots of an affine Lie algebra and ni are Kac marks. For
  crystallographic Coxeter groups (ADE series), these theories are well understood.

  [L3 --- Model-derived] Fring and Korff (2005) extended the construction to
  non-crystallographic Coxeter groups H2, H3, and H4 via a reduction (folding)
  procedure     (Nucl.   Phys.     B    729     (2005)   361,    ScienceDirect      DOI:
  10.1016/j.nuclphysb.2005.08.044). Their central result is that the classical mass
  spectra of the H2, H3, and H4 Toda theories contain particle pairs whose relative
  masses differ by exactly $\varphi$, as a direct consequence of the embedding structure. The
  paper also computes the one-loop corrections to mass renormalization, which is directly
  relevant to Work Package 2.

  The Coxeter group H2 is the symmetry group of the regular pentagon; H3 is the
  icosahedral group (of order 120); H4 is the symmetry group of the 600-cell in 4D (of
  order 14400). All have $\varphi$ as an eigenvalue of their Cartan-analogue matrices by
  construction.

  This provides a second, independent field-theoretic pedigree for exact $\varphi$ occurrence ---
  distinct from the E8 Ising model in that it arises from a geometric reduction rather than
  from E8 representation theory directly.

  4.4 Evidence Hierarchy and Mechanism-by-Mechanism Falsification
  Table
  The following table codifies, for each mechanism candidate, what outcomes would
  validate, downgrade, or falsify the claim. Rung assignments follow the derivation ladder
  of Section 3.5. This table governs the interpretation of Work Package results.